% !TeX root = ../../Book1.tex
% 纲要标题：### 13.3 典型局部化
% 纲要位置：计划纲要.md，第 484 行。

\section{典型局部化}
\label{b1:sec:1303}

分母集的选择决定哪些差别会被保留。本节考察两种常用选择：避开一个素理想的全部元素，以及一个固定元素的各次幂。前者把注意力集中于一个素理想，后者只要求一个元素可逆。整数环和多项式环的具体计算会说明，这些构造虽然采用同一套分式运算，却有不同的用途。

% 纲要内容（仅作写作计划，尚非教材正文）：
%
% * 素理想处的局部环
% * 单个元素处的局部化
% * 整数的局部化与 Laurent 多项式环
%

\subsection{素理想处的局部环}
\label{b1:subsec:1303:01}

设 \(\mathfrak p\) 为交换环 \(R\) 的素理想。因为 \(\mathfrak p\neq R\)，\(1\notin\mathfrak p\)；又因为
\(ab\in\mathfrak p\Longrightarrow a\in\mathfrak p\) 或 \(b\in\mathfrak p\)，其补集 \(R\setminus\mathfrak p\) 对乘法封闭。于是可以定义
\[
R_{\mathfrak p}\coloneqq(R\setminus\mathfrak p)^{-1}R.
\]
\symidx{\(R_{\mathfrak p}\)：环 \(R\) 在素理想 \(\mathfrak p\) 处的局部化}
它称为 \(R\) 在 \(\mathfrak p\) 处的局部化。由于 \(0\in\mathfrak p\)，分母集不含零，所以\cref{b1:prop:localization-zero-unit}保证 \(R_{\mathfrak p}\) 不是零环。

\begin{definition}\label{b1:def:local-ring}
设 \(A\) 为非零交换含幺环。如果 \(A\) 恰有一个极大理想 \(\mathfrak m\)，就称 \(A\) 为\term[jubu-huan]{局部环}[local ring]，并常以 \((A,\mathfrak m)\) 同时注明环和它的唯一极大理想。此时商环 \(A/\mathfrak m\) 是域，称为 \(A\) 的\term[shengyu-yu]{剩余域}[residue field]。
\end{definition}
这里商环确为域，是因为\cref{b1:prop:prime-max-quotient}已证明极大理想的商为域。域自身也是局部环，其唯一极大理想为零理想；“局部环”并不意味着一定存在非零非单位元。

\begin{proposition}\label{b1:prop:localization-at-prime}
设 \(R\) 为交换含幺环，\(\mathfrak p\) 为 \(R\) 的素理想。则 \(R_{\mathfrak p}=(R\setminus\mathfrak p)^{-1}R\) 是局部环，唯一极大理想为
\[
\mathfrak pR_{\mathfrak p}
=\Bigl\{\,a/s\mid a\in\mathfrak p,\ s\in R\setminus\mathfrak p\,\Bigr\}.
\]
对 \(a\in R\)、\(s\in R\setminus\mathfrak p\)，分式 \(a/s\) 是单位当且仅当 \(a\notin\mathfrak p\)。
\end{proposition}
\begin{proof}
若 \(a\notin\mathfrak p\)，则 \(a\) 也可作分母，\(s/a\) 就是 \(a/s\) 的逆。若 \(a\in\mathfrak p\) 而 \(a/s\) 为单位，\cref{b1:prop:localization-zero-unit}给出某个 \(b\in R\) 使 \(ab\notin\mathfrak p\)，与理想的吸收性矛盾。

记陈述中的集合为 \(\mathfrak n\)。两个分子在 \(\mathfrak p\) 中的分式相加时，通分后的分子仍在 \(\mathfrak p\) 中；取负及乘以任意分式也保持这一性质，所以 \(\mathfrak n\) 是理想。刚才的单位判据表明它恰为全部非单位组成的集合，因此不含 \(1\)，是真理想。任何真理想都不能含有单位，从而必包含于 \(\mathfrak n\)；这同时说明 \(\mathfrak n\) 极大，而且是唯一极大理想。它也正是由所有 \(a/1\)、\(a\in\mathfrak p\) 生成的理想，故可记为 \(\mathfrak pR_{\mathfrak p}\)。
\end{proof}

进一步可以算出这个局部环的剩余域。记 \(\bar a\) 为 \(a\) 在 \(R/\mathfrak p\) 中的像。这个商是整环，故有分式域。映射
\[
R\longrightarrow\operatorname{Frac}(R/\mathfrak p),\qquad
a\longmapsto\bar a/\bar1
\]
把每个 \(s\notin\mathfrak p\) 送到非零元素，从而送到单位。由\cref{b1:thm:localization-universal}，它唯一延拓为
\[
R_{\mathfrak p}\longrightarrow\operatorname{Frac}(R/\mathfrak p),
\qquad a/s\longmapsto\bar a/\bar s.
\]
任何非零 \(\bar s\) 都有 \(s\notin\mathfrak p\) 的代表，故此映射满射；其核由 \(\bar a=0\) 的分式组成，正是 \(\mathfrak pR_{\mathfrak p}\)。\cref{b1:thm:ring-first-iso}将商环同构到这个实际像，得到
\begin{equation}\label{b1:eq:local-residue-field}
R_{\mathfrak p}/\mathfrak pR_{\mathfrak p}
\cong\operatorname{Frac}(R/\mathfrak p).
\end{equation}
这说明素理想处局部化的剩余域一般不只是 \(R/\mathfrak p\)；当 \(\mathfrak p\) 极大时，后一商本来就是域，才无须再增加分母。

\subsection{单个元素处的局部化}
\label{b1:subsec:1303:02}

设 \(R\) 为交换含幺环，\(f\in R\)。令 \(S=\{\,f^n\mid n\geq0\,\}\)，其中 \(f^0=1\)。称局部化 \(S^{-1}R\) 为 \(R\) 在元素 \(f\) 处的局部化，记为 \(R_f\) 或 \(R[1/f]\)；其中每个元素都能写成 \(a/f^n\)，\(a\in R\)、\(n\geq0\)。
\symidx{\(R_f=R[1/f]\)：使元素 \(f\) 可逆的局部化}
两种记号都很常用；写 \(R[1/f]\) 时，\(1/f\) 表示构造中的形式逆元，不预先假定它已经存在于某个包含 \(R\) 的域中。例如 \(f\) 是零因子时，\(R\rightarrow R_f\) 可能不单射。

\begin{proposition}\label{b1:prop:single-element-localization}
设 \(R\) 为交换含幺环，\(f\in R\)，并令
\[
R_f=\{\,1,f,f^2,\ldots\,\}^{-1}R.
\]
则有保持 \(R\) 中元素像的同构
\[
R_f\cong R[t]/(ft-1),
\]
其中 \(t\) 为一个多项式变量，\(1/f\) 对应于 \(t\) 的剩余类。此外，\(R_f\) 为零环当且仅当 \(f\) 幂零。
\end{proposition}
\begin{proof}
记 \(B=R[t]/(ft-1)\)，并用上横线表示剩余类。在 \(B\) 中有 \(\bar f\bar t=1\)，故自然映射 \(R\rightarrow B\) 把每个 \(f^n\) 送到单位。\cref{b1:thm:localization-universal}给出唯一同态
\[
\alpha:R_f\longrightarrow B,\qquad a/f^n\longmapsto\bar a\bar t^{\,n}.
\]
另一方面，多项式求值
\[
R[t]\longrightarrow R_f,\qquad
\sum_{j=0}^m a_jt^j\longmapsto\sum_{j=0}^m\frac{a_j}{f^j}
\]
保持加法、乘法与单位，并把 \(ft-1\) 送到零，因此诱导同态 \(\beta:B\rightarrow R_f\)。复合 \(\beta\alpha\) 固定每个 \(a/f^n\)，所以是恒等映射；复合 \(\alpha\beta\) 固定全部 \(\bar a\) 和 \(\bar t\)，故固定它们组成的每个多项式，即固定 \(B\) 的全部元素。两者互逆。

由\cref{b1:prop:localization-zero-unit}，\(R_f=0\) 当且仅当其分母集中含零，即某个 \(f^n=0\)。这恰为 \(f\) 幂零；若 \(n=0\) 已使 \(1=0\)，则 \(R\) 本来就是零环，也仍满足这一结论。
\end{proof}

例如对 \(R=\mathbb Z/6\mathbb Z\)、\(f=\bar3\)，模 \(2\) 的映射将 \(\bar3\) 送到 \(1\)，所以延拓为 \(R_{\bar3}\rightarrow\mathbb F_2\)。其中 \(\bar2/1=0\)，每个分母 \(\bar3^n\) 的像又等于 \(1\)；因而局部化中每个元素都为 \(0\) 或 \(1\)。上述映射把两者分别送到 \(0\)、\(1\)，于是 \(R_{\bar3}\cong\mathbb F_2\)。

\subsection{整数的局部化与 Laurent 多项式环}
\label{b1:subsec:1303:03}

对一个素数 \(p\)，整数环在 \((p)\) 处的局部化为
\[
\mathbb Z_{(p)}
=\{\,a/b\in\mathbb Q\mid a,b\in\mathbb Z,\ p\nmid b\,\}.
\]
它的单位是其中分子也不被 \(p\) 整除的分式，唯一极大理想是 \(p\mathbb Z_{(p)}\)，剩余域为 \(\mathbb F_p\)。这些断言分别来自\cref{b1:prop:localization-at-prime}和\eqref{b1:eq:local-residue-field}。例如 \(2/3\) 在 \(\mathbb Z_{(5)}\) 中是单位，\(5/3\) 则不是。分母可含任意不被 \(5\) 整除的整数，不限于某个固定整数的幂。

\ProseParagraphBreak
与此不同，\(\mathbb Z[1/p]\) 只允许 \(p\) 的幂作分母：
\[
\mathbb Z[1/p]=\{\,a/p^n\in\mathbb Q\mid a\in\mathbb Z,\ n\geq0\,\}.
\]
前一个环使 \(p\) 以外的素数可逆，后一个环使 \(p\) 可逆。因而 \(\mathbb Z_{(p)}\) 和 \(\mathbb Z[1/p]\) 表示相当不同的环。更一般地，\(\mathbb Z[1/n]\) 允许 \(n\) 的所有素因子取逆；例如 \(\mathbb Z[1/6]\) 的元素可以写成 \(a/(2^r3^s)\)，其中 \(r,s\geq0\)，因为这类分母都整除足够高次的 \(6\) 的幂。

\ProseParagraphBreak
设 \(k\) 为域。在 \(k[t]\) 中使 \(t\) 可逆，得到\term[Laurent-duoxiangshi-huan]{Laurent 多项式环}[Laurent polynomial ring]
\[
k[t,t^{-1}]\coloneqq k[t]_t.
\]
\symidx{\(k[t,t^{-1}]\)：Laurent 多项式环}
其中每个元素有形式
\[
\sum_{j=m}^n a_jt^j,\qquad m,n\in\mathbb Z,\quad a_j\in k,
\]
仅有有限多个非零系数；不允许任意无限级数。因为分式 \(f(t)/t^r\) 可把 \(f\) 的每个指数减去 \(r\)，它们恰好给出这些有限和。若两个有限和相等，乘以足够大的 \(t\) 的幂后，就成为 \(k[t]\) 中的多项式等式。自然映射 \(k[t]\rightarrow k[t]_t\) 单射，所以同次幂的系数分别相等，有限和表达由此唯一。

\begin{proposition}\label{b1:prop:laurent-units}
设 \(k\) 为域。Laurent 多项式环 \(k[t,t^{-1}]\) 的单位恰为 \(ct^r\)，其中 \(c\in k^\times\)、\(r\in\mathbb Z\)。
\end{proposition}
\begin{proof}
这样的元素有逆 \(c^{-1}t^{-r}\)。反过来，设非零 Laurent 多项式 \(u\) 的最低、最高非零次数分别为 \(m,n\)，而 \(v\) 的为 \(r,s\)。因为 \(k\) 是域，乘积的最低和最高项系数都不为零，故 \(uv\) 的两个端点次数分别为 \(m+r,n+s\)。若 \(uv=1\)，则二者均为零，因而
\[
(n-m)+(s-r)=0.
\]
两个括号都是非负整数，只能各为零。所以 \(m=n\)，\(u\) 只有一个非零项，必为 \(ct^m\)。
\end{proof}
例如 \(t\) 与 \(t^{-1}\) 都是单位，而 \(t+1\) 不是。局部化只保证指定分母及其必要因子可逆，并不会自动成为分式域 \(k(t)\)。
